%=======================================================================
%  LAST-MINUTE REVISION
%  Master-only. Every count here is the one the chapter itself prints (listed
%  by src\tiers.py), and each of those was checked against the Detailed PYQ by
%  src\check.py. Nothing in this section is new material: every line points
%  back into a chapter. House style: no em dashes.
%=======================================================================
\cleardoublepage
\section{Last-Minute Revision}

{\color{sub}\footnotesize Built from the frequency analysis of all 29 papers (2065--2082 BS),
ME 708. Nothing here is new: every line points back into a chapter.}

\hr

\T{R.1 Must Know \mk{(every TOP-tier topic: these carry the paper)}}

{\color{sub}\footnotesize A topic is TOP tier when \textbf{8 or more} of the 29 papers ask it.
There are 27 of them. Every chapter has at least one; Chapter 1 has twelve.}

\vspace{2pt}
{\small
\begin{tabularx}{\textwidth}{@{}L{0.6cm}L{6.4cm}L{1.1cm}X@{}}
\toprule
\hrow \thead{Ch} & \thead{Topic} & \thead{Papers} & \thead{The one thing to be able to write} \\
2 & \textbf{Personnel management: definition, functions, role} & \textbf{17} & Flippo's definition; managerial + 7 operative functions \\
5 & \textbf{MIS: definition, need, functions, importance} & \textbf{17} & Need by manager type (production, marketing, personnel) \\
1 & \textbf{Organization: definition, necessity, importance in society} & \textbf{16} & Mooney-Reiley; 8 reasons for necessity; role in society \\
3 & \textbf{Motivation: definition, human needs, types} & \textbf{16} & Need $\rightarrow$ tension $\rightarrow$ drive $\rightarrow$ behaviour $\rightarrow$ satisfaction \\
1 & \textbf{Functions of management} & \textbf{15} & POSDCoRB list; Fayol's five for ``basic'' functions \\
1 & \textbf{Marketing: definition, concept, importance} & \textbf{14} & Stanton; 4 P's; selling vs marketing concept; digital era \\
2 & \textbf{Recruitment and selection process} & \textbf{13} & 11 steps from job analysis to induction; recruitment vs selection \\
2 & \textbf{Factors affecting wage / salary structure} & \textbf{12} & 10 factors with a Nepali example each \\
3 & \textbf{Leadership: definition, need, qualities} & \textbf{12} & 12 qualities; leader vs regular employee \\
1 & Formal vs informal organization & 11 & 10-row difference table; ``do we need informal?'' \\
1 & Scientific management (Taylor) & 11 & 5 principles, 7 techniques, pig-iron numbers \\
1 & Behavioural approach (Mayo, Hawthorne) & 11 & 4 experiments; 3-approach comparison table \\
3 & Leadership styles; which style where & 11 & Autocratic, democratic, free-rein; situational answer \\
1 & Managerial skills by level & 10 & Katz: technical low, human all, conceptual top; the figure \\
1 & Joint stock company & 10 & Features, private vs public (101 / 7 min), OCR registration \\
1 & Purchasing: 5 R's, functions & 10 & 5 R's; 13 functions of purchasing department \\
2 & Job analysis, description, specification & 10 & JD = the job, JS = the person; lecturer example \\
2 & Incentives & 10 & Financial / non-financial; individual / group \\
3 & Herzberg two-factor theory & 10 & Hygiene vs motivators lists; ``motivation comes from the job'' \\
1 & Principles of organization & 9 & 15 principles with an example each \\
1 & Levels of management & 9 & Top / middle / lower duties; time-on-functions figure \\
1 & Fayol's 14 principles & 9 & All 14 in order; applicable today \\
3 & Techniques of motivation & 9 & 9 techniques; attitude / group / executive \\
3 & Maslow's hierarchy & 9 & 5 levels + what the firm gives at each; pyramid \\
2 & Manpower planning & 8 & 7 steps; importance \\
3 & Entrepreneurship and entrepreneur's characteristics & 8 & Schumpeter; 12 characteristics \\
4 & Objectives of case study & 8 & 7 objectives; value \\
\bottomrule
\end{tabularx}}

\par\penalty0\vspace{4pt}

%=======================================================================
\T{R.2 Must Memorize \mk{(names, numbers and lists that cannot be improvised)}}

\sbs{%
\lead{Ch 1: Introduction}
\begin{itemize}
  \item \textbf{Mooney and Reiley}: organization = form of every human association for a common
    purpose. \textbf{Barnard}: formal = consciously coordinated activities; informal =
    without conscious common purpose.
  \item \textbf{Katz}: technical, human, conceptual. \textbf{Mintzberg}: interpersonal,
    informational, decisional roles.
  \item \textbf{7 models}: hierarchical, task-oriented, allocational, transactional, team
    effort, knowledge-oriented, goal-concentrated.
  \item \textbf{Taylor} (1856--1915), \textit{Principles of Scientific Management} 1911;
    pig iron 12.5 $\rightarrow$ 47.5 t/day, \$1.15 $\rightarrow$ \$1.85.
  \item \textbf{Fayol} (1841--1925), 1916; 14 principles; POCCC.
  \item \textbf{Mayo}, Hawthorne 1924--32: illumination, relay room, interviews (21{,}000),
    bank wiring (14 men).
  \item \textbf{Drucker} MBO 1954. \textbf{Graicunas}: $n(2^{n-1}+n-1)$, 6 $\rightarrow$ 222.
  \item \textbf{Nepal law}: Companies Act 2063 (private 1--101, public min 7, Rs 1 crore);
    Cooperative Act 2074 (min 30); Labour Act 2074.
  \item \textbf{5 R's}: quality, quantity, price, time, source.
\end{itemize}
\lead{Ch 2: Personnel Management}
\begin{itemize}
  \item \textbf{Flippo}: procurement, development, compensation, integration, maintenance.
  \item Wage levels: minimum, fair, living. Nominal vs real wage.
  \item Job evaluation: ranking, classification, factor comparison, point rating; ``rate the
    job, not the man''.
  \item Incentive plans: Taylor (two rates), Halsey (50\% of time saved), Rowan (peaks at
    $T_a = T_s/2$), Emerson (bonus from 66.67\%), Scanlon (labour / sales), Priestman
    (output / worker).
\end{itemize}}{%
\lead{Ch 3: Motivation, Leadership, Entrepreneurship}
\begin{itemize}
  \item \textbf{Maslow} 1943: physiological, safety, social, esteem, self-actualization.
  \item \textbf{ERG} (Alderfer 1969): existence, relatedness, growth; progression +
    frustration-regression.
  \item \textbf{McClelland} 1961: achievement, power, affiliation.
  \item \textbf{McGregor} 1960: X (coerce) vs Y (involve).
  \item \textbf{Herzberg} 1959, 200 engineers and accountants: hygiene (policy, supervision,
    pay, security, conditions, relations) vs motivators (achievement, recognition, work
    itself, responsibility, advancement, growth).
  \item \textbf{Vroom} 1964: $M = E \times I \times V$.
  \item \textbf{Blake and Mouton} 1964: (1,1) impoverished, (1,9) country club, (9,1)
    authority, (5,5) middle, (9,9) team.
  \item Leadership approaches: trait, behavioural, contingency (Fiedler, path-goal,
    Hersey-Blanchard), integrated.
  \item Power: legitimate, reward, coercive, expert, referent.
\end{itemize}
\lead{Ch 4: Case Studies}
\begin{itemize}
  \item \textbf{5 phases}: study, collect facts, define problem, individual decisions, learn.
  \item \textbf{6 steps}: state problem, collect and analyse data, tentative solutions,
    recommend, write report, checklist.
  \item Types: illustrative, exploratory, cumulative, critical instance.
\end{itemize}
\lead{Ch 5: MIS}
\begin{itemize}
  \item Information qualities: accurate, timely, complete, concise, relevant.
  \item \textbf{TPS} (operational) $\rightarrow$ \textbf{MIS} (middle, reports)
    $\rightarrow$ \textbf{DSS} (what-if) $\rightarrow$ \textbf{ESS} (top, projections).
  \item Information needs: operational hourly/daily, tactical weekly/monthly, strategic
    quarterly/yearly.
\end{itemize}}

\par\penalty0\vspace{4pt}

%=======================================================================
\T{R.3 Must Practice \mk{(opinion and applied questions: rehearse a structure, not a text)}}

{\color{sub}\footnotesize O\&M has no numerical in 29 papers. What it has instead is a steady
supply of ``in your opinion'' and ``prepare / recommend'' questions, worth full marks to a
student who takes a side and argues it in points.}

\vspace{2pt}
{\small
\begin{tabularx}{\textwidth}{@{}L{5.6cm}L{1.6cm}X@{}}
\toprule
\hrow \thead{Question} & \thead{Papers} & \thead{Take this side, in this order} \\
Which management theory suits Nepal / is most practical today? & 80 Ba, 74 Ash & Contingency; one reason per sector (factory, public enterprise, IT) \\
Structure for a (temporary) engineering project & 73 Ch, 72 Ch & Project / matrix; why not line, functional, committee \\
CEO of a software company: ownership and structure & 74 Ch & Private limited; matrix, draw it \\
Leadership style for an engineering project / industry / Nepal / start-ups & 82 Ba, 80 Ba, 78 Bh, 74 Ash, 73 Ch & Democratic base, autocratic in crises, situational \\
Can Maslow explain Devkota? & 69 Ch & Only partly; ERG and intrinsic motivation \\
Falling from grace while climbing Maslow's pyramid & 78 Bh & Esteem without ethics, frustration-regression, remedies \\
Job description and specification of a lecturer & 81 Ba & JD list, JS list \\
Which job analysis method suits Nepal? & 82 Bh & Interview + observation; why not questionnaire \\
Salary and benefits you expect & 82 Ba, 75 Ch & Market rate + allowances + PF / SSF + growth \\
HR manager vs talent poaching & 79 Ba & Pay, ESOP, career path, culture, succession \\
Prepare a case study of NEA / Nepal Telecom & 72 Ch & Report structure; planning horizon, leadership, motivation, HRD \\
Recommendations to your college & 74 Ch & Academic, infrastructure, support, research, management, image \\
Would you start a startup after graduation? & 74 Ch & Either, justified; then Nepal's risks \\
Why promote entrepreneurship in Nepal & 82 Bh, 81 Ba, 75 Ash & Migration, remittance, trade deficit, unused resources \\
\bottomrule
\end{tabularx}}

\lead{Figures to be able to draw in two minutes}
\begin{itemize}
  \item Line, functional, line \& staff, matrix organization charts ($\S$1.4, $\S$1.5).
  \item Levels pyramid; skills-by-level bars ($\S$1.2). Organization as an open system ($\S$1.1).
  \item Maslow's pyramid; Vroom's E-I-V chain; managerial grid ($\S$3.2, $\S$3.3).
  \item Functional-areas pyramid of information systems; TPS-MIS-DSS-ESS cycle ($\S$5.3).
\end{itemize}

\par\penalty0\vspace{4pt}

%=======================================================================
\T{R.4 Most Repeated PYQs \mk{(near-identical wording, set again and again)}}

\begin{enumerate}
  \item \textbf{Define personnel management and explain its functions.} 17 papers. $\S$2.1.
  \item \textbf{What are the factors affecting wage / salary structure? What is an
    incentive?} 12 + 10 papers, often in one question (76 Ch, 78 Bh, 74 Ash, 72 Ka, 70 Asa).
    $\S$2.7.
  \item \textbf{Define marketing and explain its importance in an organization.} 14 papers.
    $\S$1.6.
  \item \textbf{Explain the functions of the purchasing department.} 10 papers. $\S$1.6.
  \item \textbf{Principles of organization} with a rider (necessity, formal / informal).
    9 papers. $\S$1.1.
  \item \textbf{Explain Herzberg's hygiene theory} / \textbf{Maslow's hierarchy} /
    \textbf{Vroom's expectancy theory}, each after ``define motivation''. $\S$3.2.
  \item \textbf{Qualities of a good leader.} 12 papers. $\S$3.3.
  \item \textbf{Blake and Mouton's managerial grid.} 6 papers, 3 of them a full 8 marks. $\S$3.3.
  \item \textbf{Define entrepreneurship; its characteristics / steps for a small scale unit.}
    $\S$3.4.
  \item \textbf{Objectives of case study} (or \textbf{steps / phases}), often as a 4-mark
    short note. $\S$4.1, $\S$4.2.
  \item \textbf{Define MIS. Explain information support for functional areas of management.}
    $\S$5.1, $\S$5.3.
\end{enumerate}

\par\penalty0\vspace{4pt}

%=======================================================================
\T{R.5 One-Day Revision Checklist}

\sbs{%
\lead{Morning (Ch 1, about 40\% of the marks)}
\begin{itemize}
  \item[$\square$] Organization: definition, necessity, principles, formal vs informal.
  \item[$\square$] Management: functions, levels, skills (figure), models.
  \item[$\square$] Taylor, Fayol, Mayo + comparison table; contingency.
  \item[$\square$] Forms of ownership table; JSC formation in Nepal; cooperatives.
  \item[$\square$] Line / functional / line \& staff / committee / matrix + charts.
  \item[$\square$] Purchasing 5 R's, functions, methods; marketing importance, functions.
\end{itemize}
\lead{Midday (Ch 2)}
\begin{itemize}
  \item[$\square$] Functions of personnel management; personnel policy contents.
  \item[$\square$] Manpower planning steps; recruitment and selection steps.
  \item[$\square$] Job analysis, JD, JS; job evaluation vs merit rating; appraisal methods.
  \item[$\square$] Wage factors; incentive types; one incentive formula.
\end{itemize}}{%
\lead{Afternoon (Ch 3)}
\begin{itemize}
  \item[$\square$] Motivation definition, techniques.
  \item[$\square$] Maslow, ERG, Herzberg, X-Y, Vroom: one figure or table each.
  \item[$\square$] Leadership qualities, styles, grid, approaches, leader vs manager.
  \item[$\square$] Entrepreneur characteristics, small-scale steps, why promote in Nepal.
\end{itemize}
\lead{Evening (Ch 4, Ch 5)}
\begin{itemize}
  \item[$\square$] Case study objectives, 5 phases, 6 steps.
  \item[$\square$] MIS definition, need, importance; data vs information.
  \item[$\square$] TPS / MIS / DSS / ESS table; functional-areas pyramid.
  \item[$\square$] Computers and MIS; database; organizing IS.
\end{itemize}
\lead{Night before}
\begin{itemize}
  \item[$\square$] Read R.1 and R.2 only. Rehearse two opinion answers from R.3.
  \item[$\square$] Every answer: definition line, numbered points in bold, one Nepali example,
    one figure where it fits.
\end{itemize}}
